\section{Soft response of the cubic and quartic curvature terms}
\label{sec:ex-soft-curvature-receipts}

This section supplies the order-two and order-three response estimates used
by the selected one-loop route. The argument has two parts. First, the
absolute path budgets give common analytic pictures for the cubic and
quartic terms, uniformly over the volume sequence. Second, a finite
gauge-variation identity removes the unwanted low-momentum coefficients.
The latter step is needed: analyticity and an absolute activity bound alone
do not imply the extra powers of momentum.

\subsection{Summable local atoms and their coefficient pictures}

An atom of degree $r$ consists of $r$ transported curvature slots and a
finite coefficient array $c_\alpha$, where $\alpha$ records one color
coordinate for each slot. Write $d(a)$ for the sum of the slot lengths and
\[
   b_\sigma(a)=\sum_\alpha |c_\alpha|e^{\sigma d(a)}.
\]
The corresponding local function is the polynomial in those $r$ slot
values. For a sequence $a_i$, its absolute activity budget is
\[
   \sum_i b_\sigma(a_i)\leq A.
\]
Gauge invariance means that simultaneous adjoint rotation of all slot
arguments leaves each polynomial unchanged. These are precisely the two
properties of the cubic and quartic streams used below.

\begin{proposition}[Uniform coefficient picture for an atom stream]
\label{prop:ex-atom-stream-picture}
Fix $r,t\in\mathbb N$, $r\leq4$, and $\sigma>0$. There are constants
$\eta_{r,t}>0$ and $C_{r,t}<\infty$, depending only on $G,r,t,\sigma$,
such that the following holds. If $\sum_i b_\sigma(a_i)\leq A$, then the
translation sum of the local atom functions has an exact order-$(t+1)$
Fourier coefficient picture $P_{r,t}$. Its components are holomorphic on
the independent momentum strip of width $\eta_{r,t}$, have uniform bound
\[
       \|P_{r,t}\|_{\eta_{r,t}}\leq A C_{r,t},
\]
and represent the finite-volume Taylor coefficient on every torus. At
nonzero total momentum the coefficient vanishes; at zero total momentum it
is the torus cardinality times the picture value. The bound is independent
of the torus side and of the number of atoms.

If the atoms are gauge invariant, the same conclusion holds for their
orientation-averaged action, with the price multiplied by
\[
       \exp((t+1)^2\eta_{r,t}).
\]
\end{proposition}

\begin{proof}
The one-slot estimate in the proof of
Proposition~\ref{prop:ex-row-quadratic-price} gives, on one common complex
ball about the physical identity, a bound
\[
       |s_{\nu,\alpha_\nu}(z)|
       \leq S_G e^{(\sigma/2)\ell_\nu}.
\]
It was obtained from the complexified skew-color projection of the forward
slot string minus its reverse; thus it applies to the real color
coordinates and their complex-linear extension. Multiplying the $r$ slot
bounds and summing the coefficient magnitudes gives
\[
       |a_i^{\mathbb C}(z)|
       \leq S_G^r b_{\sigma/2}(a_i).
\]
The complex extension is entire in the finitely many link variables in the
atom's stencil. Cauchy's derivative estimate on nested balls of radii
$R,R/2,\ldots,R/2^{t+1}$ bounds the normalized order-$(t+1)$ derivative by
\[
 b_{\sigma/2}(a_i)S_G^r
       \frac{\operatorname{circleTax}(R,t+1)}{(t+1)!}.
\]
Here $\operatorname{circleTax}(R,0)=1$ and
$\operatorname{circleTax}(R,k+1)=\operatorname{circleTax}(R,k)/(R/2^{k+1})$;
the factorial is the Taylor coefficient normalization.

The local coefficient picture inserts one Fourier character for each of
the $t+1$ differentiated slots. A slot support lies in the four-dimensional
box of radius at most $d(a_i)$. On a momentum strip of width $\eta$, the
imaginary part of each plotted momentum is bounded coordinatewise by
$(t+1)\eta$. Hence the product of the $t+1$ characters is bounded by
\[
       \exp\bigl(4(t+1)^2\eta d(a_i)\bigr).
\]
Choose $\eta=\sigma/(8(t+1)^2)$, so
$4(t+1)^2\eta=\sigma/2$. The resulting term price is bounded by
\[
 b_{\sigma/2}(a_i) S_G^r
 \frac{\operatorname{circleTax}(R,t+1)}{(t+1)!}
 e^{(\sigma/2)d(a_i)}
 = C_{r,t} b_\sigma(a_i),
\]
where $C_{r,t}=S_G^r\operatorname{circleTax}(R,t+1)/(t+1)!$.
Thus the coefficient pictures form a normally convergent series on every
smaller strip. The same majorant on the common complex link ball permits
termwise differentiation of the atom series before extracting any
momentum coefficient. This proves analyticity and the claimed uniform
bound.

For the finite-volume identity, translate the local atom over all anchors.
The sum of its characters is the finite-torus character sum: it is zero
unless the total momentum is zero, and equals the number of sites when the
total momentum is zero. The local picture is defined on a cover box large
enough to contain the stencil, then pulled back by the wrapping map. This
keeps the identity exact even on small tori where different cover points
wrap to the same link. Absolute summability just proved permits exchanging
the atom sum, derivative, and finite anchor sum.

For an orientation $s$, the picture transforms by multiplying the $t+1$
reflected-character fees and evaluating the original picture at the
inverse-permuted direction labels and transformed independent momenta.
Each fee has strip price at most $e^{(t+1)\eta}$, because its plotted
momentum has coordinatewise imaginary parts bounded by $(t+1)\eta$.
The product therefore costs at most $e^{(t+1)^2\eta}$. The normalized
finite orientation average preserves this bound by the triangle
inequality, and its coefficient identity is the average of the exact
finite-volume identities. This proves the final assertion.
\end{proof}

\subsection{The finite gauge-variation identity}

A coefficient picture of order $t$ records the normalized $(t+1)$st
Taylor coefficient. For an external color word $v$, color run $a$, and
independent momentum tuple $p$, its one-leg contraction at leg $j$ is
\[
 (\mathcal C_jP)_{v,a}(p)=
   \sum_{\mu\in\mathcal A}q(p_j,\mu)
       P_{v,\operatorname{update}(a,j,\mu)}(p),
 \qquad q(p_j,\mu)=e^{ip_{j,\mu}}-1,
\]
where $\mathcal A$ is the finite set of spatial axes. The color word
$v$ is held fixed; the sum is over the spatial difference axis.

\begin{lemma}[Ward partition identity]
\label{lem:ex-ward-partition}
Let $m_k\to\infty$, and let $F_k$ be real analytic, gauge invariant
actions with $DF_k(0)=0$. Suppose an order-$t$ picture $P$ is uniformly
bounded on a momentum strip and records the $(t+1)$st coefficient of
$F_k$ exactly on every torus. Suppose further that for each
$b=2,\ldots,t$ there is a uniformly bounded order-$(b-1)$ picture
recording the $b$th coefficient, with each of its entries of ray order
at least $s$. Then each contraction $\mathcal C_jP$ has ray order at
least $s$.
\end{lemma}

\begin{proof}
Differentiate the composition of $F_k$ with the infinitesimal gauge
change $t+1$ times. The chain rule is the finite sum over ordered set
partitions of the $t+1$ slots. Gauge invariance makes this sum zero:
the pure-gauge bundle has no component in the distinguished gauge
coordinate. The all-singleton partition is, by the exact coefficient
record, $-|\Lambda_{m_k}|(t+1)!\,\mathcal C_jP$. The factor
$q(p_j,\mu)$ is precisely the Fourier multiplier of the one-link
difference in the distinguished slot. The one-block term vanishes since
$DF_k(0)=0$.

For a partition with $b$ blocks, $2\leq b\leq t$, the lower recorded
picture has order $b-1$ because it records the $b$th derivative. Its
arguments are the sums of the external momenta in the corresponding
blocks. The chain-rule contribution is a finite sum of products of this
picture entry and continuous block-derivative coefficients; it therefore
has ray order at least $s$. Thus on the finite period mesh the
contraction is a finite sum of such lower terms divided by the positive
factor $|\Lambda_{m_k}|(t+1)!$. Every real momentum tuple is a limit of
zero-total-momentum mesh tuples: round the $t$ independent momenta to the
period mesh, and set the remaining momentum (the first, distinguished
leg) to their negative sum, retaining the integer winding in the momentum
lift. Exact recording gives the same factor $|\Lambda_{m_k}|$ in the top
and each lower coefficient. Cancel this common volume factor in the
finite-mesh identity before passing to the limit. The lower terms are
continuous, so the identity passes to every real tuple. Their finite sum
has ray order $s$, proving the claim.
\end{proof}

\begin{lemma}[Recovering powers from the contracted kernel]
\label{lem:ex-ward-rays}
Assume the picture entries are analytic at zero. For every fixed color
word and run with at least one independent momentum variable, contractions
of ray order at least two imply that each entry has ray order at least
one. If the contractions have ray order at least three and there are at
least two independent momentum blocks (equivalently, at least three
external legs), then each entry has ray order at least two.
\end{lemma}

\begin{proof}
The first-order expansion $q(\lambda p_j,\mu)=i\lambda p_{j,\mu}
+O(\lambda^2)$ shows, after dividing a contraction of order two by
$\lambda$, that
\[
 \sum_{\mu\in\mathcal A}p_{j,\mu}
 P_{v,\operatorname{update}(a,j,\mu)}(0)=0.
\]
Choose $p$ supported at the unique axis which updates $a_j$ back to
itself. Then $P_{v,a}(0)=0$. Analyticity gives ray order one.

For a contraction of order three, the next coefficient, using that
constant terms vanish, is
\[
 \sum_{\mu\in\mathcal A}p_{j,\mu}
 D P_{v,\operatorname{update}(a,j,\mu)}(0)[p]=0. \tag{WR}
\]
Fix leg $j$ and a vector $w$ whose $j$th block is zero. Let $q$ be the
coordinate unit vector on leg $j$ at the axis that updates $a_j$ back
to itself. Evaluate (WR) at $q+w$ and subtract its value at $q$. Only
the selected axis remains, giving
$DP_{v,a}(0)[w]=0$. The equation at $q$ also gives
$DP_{v,a}(0)[q]=0$. Thus the derivative vanishes on all vectors
transverse to leg $j$. With two distinct legs, split any vector into a
part supported on the first and a part with that leg zero; apply the
transverse conclusion using the other leg for the first part and the
first leg for the second. The full differential vanishes. Together with
the vanishing constant term, this gives ray order two.
\end{proof}

\subsection{Cubic and quartic terms}

Let $A^{(3)}_i$ and $A^{(4)}_i$ be gauge-invariant atom streams of degrees
three and four, each with activity budget at most $A$ at width $\sigma$.
Write $T_r$ for their finite signed-orientation-averaged translation sums.

\begin{proposition}[Two soft pictures for the curvature sum]
\label{prop:ex-extra-curves}
If $m_k\to\infty$, there are $\eta_*>0$ and $K_*<\infty$, depending
only on $G,\sigma$, such that
\[
        F_k=T_3(m_k)+T_4(m_k)
\]
has exact recorded pictures $P_2$ of order $2$ and $P_3$ of order $3$ with
uniform price at most $A K_*$. The first picture has vanishing order $2$;
the second has vanishing order $1$. The functions $F_k$ are analytic,
gauge invariant, and vanish to order at least three at zero.
\end{proposition}

\begin{proof}
Apply Proposition~\ref{prop:ex-atom-stream-picture} to each of the four
choices
\[
       (r,t)=(3,2),(3,3),(4,2),(4,3).
\]
For each fixed $t$, add the degree-three and degree-four pictures. Their
recording identities add exactly, and their prices add. Choose the minimum
of their four positive widths and the sum of their four finite price
constants; this yields one $\eta_*$ and one $K_*$, uniform over the stream
and all periods.

The individual atom series vanish to at least their degree at zero, because
each transported curvature slot vanishes there. Absolute convergence of
the differentiated series preserves these vanishing orders, so $F_k$ has
order at least three. Gauge invariance is preserved by the finite
orientation average. The zero picture records its second Taylor
coefficient, and the zero first derivative is automatic. Thus the
order-two picture supplies the lower-order coefficient needed for the
order-three Ward argument.

Apply Lemma~\ref{lem:ex-ward-partition} with $t=2$ to the order-two
picture. The zero lower coefficient and the order-three vanishing of
$F_k$ give contractions of order three. Lemma~\ref{lem:ex-ward-rays}
then upgrades the order-two picture to vanishing order two. For the
order-three picture, apply Lemma~\ref{lem:ex-ward-partition} with $t=3$;
the zero second coefficient and the order-two picture just obtained supply
the lower-order data at block lengths two and three. The contractions
vanish to order two, and Lemma~\ref{lem:ex-ward-rays} gives vanishing
order one for the order-three picture. These are the two claimed Ray
orders.
\end{proof}

\subsection{The curvature stream and the Wilson subtraction}

For a path family $R$, write
\[
  C_k=\sum_{i\geq0}\operatorname{turn}(R^{(3)}_i)(m_k)
      +\sum_{i\geq0}\operatorname{turn}(R^{(4)}_i)(m_k),
  \qquad
  I_k=\operatorname{intended}(p,R;m_k),\quad W_k=\operatorname{Wilson}(m_k).
\]
The path bounds provide the cubic and quartic activity budgets at most
$A$ and their gauge invariance. The curvature part satisfies
$I_k=J_k+C_k$, where $J_k$ is the quadratic path action.

\begin{proposition}[Curvature stream pictures]
\label{prop:ex-veneer}
Assume the cubic and quartic budgets and invariance just stated. There are
constants $\eta_c>0$, $K_c<\infty$ such that for $k=1,2,3$ the curvature
sequence $C_k$ has a priced, exactly recorded picture of order $k$,
vanishing order $4-k$, with price at most $A K_c$.
\end{proposition}

\begin{proof}
Proposition~\ref{prop:ex-extra-curves} supplies the order-two and
order-three pictures. For order two its vanishing order is two, as
required. For order three its vanishing order is one. The order-one
picture is identically zero: $C_k$ vanishes to third order, so its second
Taylor coefficient is zero. The recorded error is therefore zero, its
price is zero, and it has any positive vanishing order. Taking the common
minimum width and the maximum of the two prices proves the assertion.
\end{proof}

Now impose the selected metric hypotheses: a bed $B$ is attached to $R$;
its metric row is bounded on a strip of width $\sigma$ and price $C$;
the metric floor bound holds at every torus side and the sworn condition
holds at each selected side $m_k$. The following estimates are the only uses of these
assumptions. The metric-slot picture has price
$b_0=\operatorname{bedCharge}(\sigma,C)$ and records the quadratic
coefficient of $J_k$. The unit leader has price
$b_1=\operatorname{bedCharge}(\sigma,1)$ and records the Wilson
coefficient. The floor-start condition gives the zero-momentum normalization
$B(0)=I$. The sworn identity and the attached-path Hessian identity
identify the actual quadratic coefficient of $J_k$ with the metric-slot
picture. In Fourier variables, the difference of the two quadratic
pictures has the form
\[
       \widehat{\operatorname{curl}}(p)^*
       (B(p)-I)\widehat{\operatorname{curl}}(p).
\]
Each curl factor is a one-link difference and is $O(|p|)$; analyticity
and $B(0)=I$ give $B(p)-I=O(|p|)$. Their product therefore has ray
order three. This is the order-one picture for $J_k-W_k$, with price
$b_0+b_1$. The higher column pictures are supplied by
Proposition~\ref{prop:ex-row-quadratic-price}; they record the third and
fourth coefficients of $J_k$. The bound
$\operatorname{bedCharge}(\sigma,D)$ is the product of the two face
counts and the elementary exponential strip bound for the two slot
factors, so it is independent of the torus side.

The row estimate of Proposition~\ref{prop:ex-row-quadratic-price}, with
its signed average, supplies priced pictures for $J_k$ at orders two and
three. Proposition~\ref{prop:ex-wilson-stream} supplies the corresponding
Wilson pictures at both orders, with price constants $X_2,X_3$ and
positive width one. Subtracting pictures preserves exact recording and
adds their prices.

\begin{proposition}[The selected Wilson-subtracted response]
\label{prop:ex-greyson}
Assume that $R$ has path bounds $(\tau,A,\Lambda)$. Its attached bed has
metric strip bound $(\sigma,C)$, the floor holds at every torus side, and
the sworn condition holds at each $m_k$. Also suppose
\[
                    0\leq p.\mathrm{direct}\leq Z.
\]
Then there are $s>0$ and $M<\infty$, uniform in the sequence and path,
such that $I_k-W_k$ has exact recorded pictures of orders $j=1,2,3$,
each priced at width $s$ and price $M$, with vanishing order $4-j$.
One valid price is
\[
 M=\operatorname{bedCharge}(\sigma,C)+\operatorname{bedCharge}(\sigma,1)
      +P_2+P_3+X_2+X_3+A K_c,
\]
where $P_2,P_3$ are the two row-action prices from
Proposition~\ref{prop:ex-row-quadratic-price}, and $K_c$ is from
Proposition~\ref{prop:ex-veneer}. One valid width is the minimum of
$\sigma$, $1$, the two row widths, and $\eta_c$.
\end{proposition}

\begin{proof}
Set $F_k=J_k-W_k$. Let $S_1$ be the metric-slot picture minus the unit
leader picture. It records the second Taylor coefficient of $F_k$, has
price $b_0+b_1$, and has vanishing order three by the floor and zero-metric
normalization. In the notation of the Ward lemma it is the lower picture
of order one for $F_k$. Let $S_2$ be the row picture for $J_k$ minus the
order-two Wilson picture, and let $S_3$ be the order-three row picture
minus the order-three Wilson picture. They exactly record the third and
fourth Taylor coefficients of $F_k$ and have prices at most $P_2+X_2$
and $P_3+X_3$.

The action $F_k$ is gauge invariant and has zero first derivative. For
$t=2$, apply Lemma~\ref{lem:ex-ward-partition} to $S_2$, with the single
lower picture $S_1$; the resulting contraction has vanishing order three,
and Lemma~\ref{lem:ex-ward-rays} gives vanishing order two for $S_2$.
For $t=3$, apply the partition lemma to $S_3$, using $S_1$ and the now
improved $S_2$ at its two lower block counts. The contraction has
vanishing order two, so the same recovery lemma gives vanishing order one
for $S_3$. This supplies the exact three records and prices required for
$F_k$.

The cubic and quartic terms $C_k$ have the order-one, order-two, and
order-three pictures of Proposition~\ref{prop:ex-veneer}. Add the
order-$j$ picture of $C_k$ to the order-$j$ picture of $F_k$ for each
$j=1,2,3$. The sum records
$F_k+C_k=I_k-W_k$ exactly. Its price is bounded by the displayed $M$,
and its vanishing order is at least $4-j$, since the order of a sum is
at least the minimum order of its summands. Taking the common minimum
strip width gives $s$.
\end{proof}

\subsection{From coefficient pictures to local response arrays}

We state the finite conversion in its concrete two-sided form. For
$j=1,2,3$, a coefficient picture $P_j$ records exactly the normalized
$(j+1)$st derivative of $U_m$ on every finite torus. It is uniformly
priced at width $s$ if each entry is holomorphic and bounded by the same
constant on the independent-momentum strip of width $s$. It has ray
order $r$ if, on every real ray through zero, each entry is
$O(|\lambda|^r)$. For an analytic entry this means every Taylor
homogeneous polynomial below degree $r$ vanishes: evaluating that
polynomial on each real direction gives zero, hence the polynomial is
identically zero. Thus ray order $r$ supplies exactly the
Taylor-vanishing hypothesis used in Lemma~\ref{lem:ex-soft-factorization}.
The required input at level $j$ is ray order $3-j$.

For $n=0,1,2$, an external color address is an $(n+2)$-tuple of
one-link color-coordinate labels, each a pair consisting of a spatial
axis and a Lie-color basis index. The multilinear test fields $w_r$
carry the corresponding link and site dependence.
The stitch address set $\mathsf{St}_n$ is the set of ordered
$(2-n)$-tuples from $\{0,\ldots,n\}\times\mathcal A$. The fixed rule
$\ell_n$ routes each key $(r,\mu)$ to its assigned response leg and axis:
for $n=0$ the two keys are routed to legs zero and one; for $n=1$ its
single key $(r,\mu)$ is routed to leg $r+1$; and for $n=2$ there are no
keys. Let $\mathsf{Shift}_n$ be the integer lattice on the $(n+1)$
independent momentum-axis coordinates. A response
array is a map
\[
 F_n:\mathsf{Color}_{n+2}\times\mathsf{St}_n\times
       \mathsf{Shift}_n\longrightarrow\mathbb C.
\]
Its width at $\eta$ is the supremum, over the finite color and stitch
addresses, of
\[
 \sum_{z\in\mathsf{Shift}_n}e^{\eta|z|_1}|F_n(b,i,z)|.
\]
The stitched Fourier picture is
\[
 \widetilde Q_n(b,p)=\sum_{i\in\mathsf{St}_n}
  \left(\prod_{r=0}^{n+1}\prod_{e\in\ell_n(i)_r}
       (e^{i\,\mathrm{plot}_r(p)\cdot e}-1)\right)
       \sum_z F_n(b,i,z)e^{ip\cdot z}.
 \tag{SR1}
\]
Here $\mathrm{plot}_r(p)$ are the $n+2$ external momenta: the first is
minus the sum of the $n+1$ independent momenta, and the remaining ones
are those independent momenta in order. For $n=0$ the two listed
differences are distributed one per leg; for $n=1$ there is one
difference, and for $n=2$ there are none.

To state the response identity, split a link fluctuation into two
species $x_+,x_-$ and write
$\mathsf D_m(x)=\operatorname{content}(\operatorname{input}(x_+)
\operatorname{input}(x_-))$. Let $v_n$ put species $-$ in the first $n$
slots and species $+$ in the final two. Let $\mathcal R_{n,m}(U_m)$ be the $(n+2)$-linear derivative of
$U_m\circ\mathsf D_m$ at zero, restricted to those species in that
order; it is a tensor whose components are indexed by external color
addresses $b$. For a list of axes $\ell=(e_1,\ldots,e_r)$, let
$\Delta_\ell=(T_{e_1}-I)\cdots(T_{e_r}-I)$ be the corresponding product
of translated one-link differences. For a test tuple $w=(w_0,\ldots,w_{n+1})$,
define
\[
 \operatorname{Local}_{m,\ell_n(i)}(F_n(b,i,\cdot))[w]
 =\sum_z F_n(b,i,z)\sum_{x\in\Lambda_m}
   \prod_{r=0}^{n+1}
   (\Delta_{\ell_n(i)_r}w_r)
      (x+\operatorname{wrap}_m(\operatorname{lane}_r(z)),b_r).
\]
Here $\operatorname{lane}_0(z)=0$, while $\operatorname{lane}_r(z)$ for
$r\geq1$ is the shift in the $(r-1)$st independent momentum coordinate.
The width bound makes the Laurent sum absolutely convergent. The response
identity required below is componentwise and multilinear:
\[
 (\mathcal R_{n,m}(U_m))_b[w]
   =\sum_{i\in\mathsf{St}_n}
       \operatorname{Local}_{m,\ell_n(i)}(F_n(b,i,\cdot))[w]
       \qquad(b\in\mathsf{Color}_{n+2}).
 \tag{SR2}
\]
\begin{lemma}[Finite difference division]
\label{lem:ex-soft-factorization}
Let $E$ be a finite set of phase coordinates. Suppose $f$ is $2\pi$-periodic in each coordinate and holomorphic,
bounded by $A$ on a strip of width $s>0$, and its Taylor expansion at zero
has no terms of degree below $q$. Then on a narrower strip it has a finite
representation
\[
 f(p)=\sum_{(a_1,\ldots,a_q)\in E^q}
        \prod_{r=1}^q(e^{ip_{a_r}}-1)\,g_{a_1,\ldots,a_q}(p),
\]
where the residual functions are holomorphic on width at least
$s/2^q$ and uniformly bounded there by $A C(E,q,s)$. For $q=0$ take the
one-term representation $g=f$.
\end{lemma}

\begin{proof}
For $J\subseteq E$, write $p^{[J]}$ for $p$ with the coordinates in $J$
set to zero. For each $a$, periodicity implies that
$f(p)-f(p^{[\{a\}]})$ vanishes whenever $p_a\in2\pi\mathbb Z$.
Consequently it is divisible, as a holomorphic periodic function, by
$d_a(p)=e^{ip_a}-1$. The quotient is removable at every zero. On a
fundamental real period cell, away from fixed neighborhoods of these
zeros the denominator is bounded below; inside those neighborhoods the
quotient is bounded by Cauchy's estimate on a narrower strip. Thus
\[
 f(p)-f(p^{[\{a\}]})=d_a(p)g_a(p)
\]
with a uniform bound for $g_a$. In the telescoping identity over $J$,
replace $f$ at each step by the function with the previously selected
coordinates staged to zero. Each corresponding quotient lowers the
vanishing order by one. Induction on $q$ applied to these residuals
therefore yields the displayed finite sum. At most $|E|^q$ terms occur;
at each division, Cauchy's estimate on the current strip of width $r$
bounds the quotient on width $r/2$; after $q$ divisions the residual
strip has width $s/2^q$. Each step multiplies the bound by a constant
depending only on $E,q,s$. The fully
staged term is zero by the order-$q$ hypothesis. This proves the result.
\end{proof}

For the local response array, the number of factors is $q=2-n$, not
$n$. If $n=1,2$, apply the factorization to every picture entry with
$q=2-n$ differences and place them on the prescribed stitch legs. For
$n=0$, two factors are required, one on each of the two response legs.
The division first gives the two factors in the single independent
momentum; the second is moved to the other leg by the exact identity
\[
 q(p,\mu)=-\operatorname{mode}(p,\operatorname{step}\mu)q(-p,\mu).
\]
The phase is absorbed into the Laurent coefficient, with its strip cost
included in the constant. Thus this construction yields precisely the
three stitch patterns: two factors for $n=0$, one for $n=1$, and none for
$n=2$.

To control the Laurent width, let $g$ be one residual, bounded by $B$
on a strip of width $\eta$. Shifting the contour in each periodic
coordinate gives $|\widehat g(\nu)|\leq B e^{-\eta|\nu|_1}$. At width
$\eta/2$,
\[
 \sum_{\nu}e^{(\eta/2)|\nu|_1}|\widehat g(\nu)|
 \leq B\left(\sum_{k\in\mathbb Z}e^{-(\eta/2)|k|}\right)^{|E|}
 =B\left(1+\frac{2}{e^{\eta/2}-1}\right)^{|E|}.
\]
Thus each residual defines a Laurent entry with a quantitative width
bound; summing the at most $|E|^q$ residuals multiplies it by that finite
factor. This is the coefficient-norm estimate used in the array
construction.

We now verify the response identity and its price using the actual
mixed-field composition. Let $D(x)$ be the local map which forms the
product of the two split fields and then takes its logarithmic link
coordinate. For a Boolean word $v$ on the external slots, the response
of $U_m$ is the corresponding mixed derivative of $U_m\circ D$ at zero.
At level $n$, the word has $n$ entries of the first species followed by
two entries of the second species; its length is $n+2$.

For a set partition $\pi$ of these $n+2$ slots into $b$ nonempty blocks,
order the blocks canonically by their least slot. The chain rule gives
the derivative of $U_m$ of order $b$ evaluated at the sums of external
momenta in the blocks, multiplied by the derivatives of $D$ on each
block. The single-block term is zero because $DU_m(0)=0$. For
$2\leq b\leq n+2$, the action coefficient has picture order $b-1$.
Its available ray order is $3-(b-1)=4-b$, which is at least
$2-n$. Thus every term in the mixed response picture has ray order
$2-n$ and is analytic on a common narrower strip.

More explicitly, for a partition $\pi$ and a basis phase assignment $w$,
the corresponding mixed coefficient is
\[
 b!\sum_{a_1,\ldots,a_b}
       \left(\prod_{r=1}^{b}c_{\pi,r}(a_r;w_r)\right)
       P_{b-1}\bigl(a_1,\ldots,a_b;
            \textstyle\sum_{j\in \pi_1}p_j,\ldots,
                      \sum_{j\in \pi_b}p_j\bigr),
\]
where $c_{\pi,r}(a_r;w_r)$ is the coordinate coefficient of the derivative
of $D$ on block $r$. The factorial is the normalization between the
symmetric derivative picture and the raw $b$th action derivative; it does
not enumerate permutations of the blocks. If the input pictures are
bounded by $A$, its absolute price is
bounded by
\[
 A b!\sum_{a_1,\ldots,a_b}
       \prod_{r=1}^{b}|c_{\pi,r}(a_r;w_r)|.
\]
Summing the absolute coefficient multiplier (the displayed sum without
the input price $A$) over partitions, color basis labels, and phase basis
labels defines the response constant $P_n$.
These coefficients are uniformly finite: the map $\mathsf D_m$ is the
same pointwise product-and-log map on each link. On a fixed identity
neighborhood in the finite-dimensional gauge group, this map is analytic
and bounded. Cauchy's estimate on a smaller ball bounds every derivative
of order at most four by a constant depending only on the group and the
order. The number of partitions and color/phase labels is finite, so
$P_n<\infty$ independently of $m$.
The common analytic strip for the mixed picture can be chosen as
\[
       s_n=\frac{s}{(n+3)^2}.
\]
Every independent external momentum has imaginary coordinates at most
$s_n$, and the missing momentum, their negative sum, has coordinates at
most $(n+1)s_n$. A block sum of at most $n+2$ external momenta therefore
has imaginary coordinates at most $(n+2)(n+1)s_n<s$. Thus the input
pictures can be evaluated at every block sum. The mixed picture is
bounded by $A P_n$ on width $s_n$ when the input picture prices are at
most $A$.

For each entry of the mixed picture $\widetilde Q_n$, finite difference
division now produces the prescribed $2-n$ stitch factors and residual
functions on width $d_n=s_n/2^{2-n}$. Shifting each contour within the
smaller strip of width $d_n/2$ gives Laurent coefficients bounded by
$B e^{-(d_n/2)|z|_1}$ when the residual supremum is $B$. At array width
$d_n/4$, their weighted sum is bounded by
\[
 B\left(1+\frac{2}{e^{d_n/4}-1}\right)^{D_n},
\]
where $D_n$ is the finite number of independent Laurent coordinates.
The number of partitions and addresses is finite, so all residuals
assemble into a response array of width at most $A C_n$ for a constant
$C_n$ depending only on the finite labels and input width.

To check exactness, evaluate the factorization on a finite torus. Each
Laurent term is absolutely summable by its weighted width. Its spatial
translation sum is a finite character sum; orthogonality makes it zero
unless total momentum is zero, in which case it equals the number of
sites. The coefficient picture recording identity then identifies the
sum with the mixed derivative of $U_m\circ D$. This is the exact
finite-volume response identity, not only equality of Taylor polynomials.

\begin{proposition}[Three-level response arrays]
\label{prop:ex-three-level-conversion}
Let $U_m$ be real analytic, with $U_m(0)=0$ and $DU_m(0)=0$. Suppose it
has exact coefficient pictures $P_j$ at orders $j=1,2,3$, whose entries
are $2\pi$-periodic in each phase coordinate, bounded on one common strip
of width $s>0$, and have ray orders at least $3-j$. Then there
are a common positive strip
width $\eta_*$ and a finite constant $C_*$ such that the response arrays
at levels $n=0,1,2$ satisfy (SR1), the exact finite-period response
identity, and width at most $C_*$ times the sum of the three input
picture prices, uniformly in $m$.
\end{proposition}

\begin{proof}
For each $n$, first form the mixed picture $\widetilde Q_n$ by the
ordered-partition chain rule. Every term is the input picture
$P_{|b|-1}$ at the block-summed momenta, multiplied by the block
coefficients. The one-block term is zero. Since $P_{|b|-1}$ has ray
order $4-|b|\geq2-n$, the mixed picture has ray order $2-n$. Apply the
finite difference lemma to $\widetilde Q_n$ with $q=2-n$ and assign the
factors to the prescribed legs; for the two-leg split use the phase
identity above. The finite character sum proves the response identity.
The Laurent estimate bounds the width. For each $n=0,1,2$, let $C_n$ be the finite price multiplier from the
ordered-partition and division construction, including the Laurent
coefficient sum. For the common input strip $s$, the mixed-picture strip at
level $n$ is $s_n=s/(n+3)^2$, and the residual strip is
$d_n=s_n/2^{2-n}$. Choose
\[
 \eta_* =\min_{0\leq n\leq2}(d_n/4),\qquad
 C_*=C_0+C_1+C_2.
\]
Then all arrays have width at most $C_*$ times the sum of the three input
picture prices. These finite constants and the positive width are
independent of the period and action.
\end{proof}

The division and response construction above outputs an array of exactly
this type. The factorization proves (SR1) entrywise. The finite-volume
character identity proves (SR2); Laurent absolute convergence justifies
the sum over shifts. Thus the exact response is supplied at each finite
period, not merely in the limiting Taylor series.

Apply the conversion to $I_m=(I_m-W_m)+W_m$. Proposition
\ref{prop:ex-greyson} supplies pictures with ray orders $3,2,1$ at
orders $1,2,3$. The Wilson pictures have ray orders $2,1,0$, so the sums
have the required orders $2,1,0$. Their exact records and prices add.
The value and first derivative of $I_m$ vanish by the quadratic
normalization and cubic vanishing of the curvature part. We obtain three
uniformly priced response arrays.
Let $M_g$ denote the displayed price in Proposition~\ref{prop:ex-greyson},
and let $P_W$ be the sum of the leader and the two higher Wilson picture
prices. The common array width is bounded by
$K=3C_*(M_g+P_W)$.

Define the mockup width property by the common Laurent bound above at each
level $n\leq2$, and define its response property to be (SR2) at every
such level. Then the theorem below supplies both properties, conventionally
called the width and exact-response requirements for the three mockups.

\begin{theorem}[Selected soft-response arrays]
\label{thm:ex-soft-response-mockups}
Fix $\sigma>0$, $M,Z,A\geq0$, $\tau>0$, and $\varepsilon\in\mathbb R$. There are
$\eta>0$ and $K<\infty$ such that for every divergent period sequence
$m_k$, clip $p$, path family $R$, attached bed $B$, and $\Lambda\geq0$
satisfying
\[
 p.\mathrm{direct}\leq Z,\quad R.\mathrm{Bounds}(p,\tau,A,\Lambda),\quad
 B\text{ attached to }R,\quad B.\mathrm{metric.strip}(\sigma,M),
\]
with $\mathrm{Floor}(B,p,n,\varepsilon,M)$ for every torus side $n$,
and the sworn metric identity at each selected side $m_k$. Put
$U_{m_k}=\operatorname{intended}(p,R;m_k)$. Then these actions admit arrays
$F_0,F_1,F_2$ of width at most $K$ satisfying (SR1) and (SR2) on every
period. Set $F_n=0$ for $n>2$. One may take $K=3C_*(M_g+P_W)$, where $M_g$ is the price displayed in
Proposition~\ref{prop:ex-greyson}, $P_W$ is the finite Wilson picture
price, and $C_*$ is the finite response-array multiplier above. The
constants $\eta,K$ depend only on the listed scalar bounds and the fixed
gauge and finite coordinate sets.
\end{theorem}

\begin{proof}
The Wilson-subtracted proposition supplies order-$j$ pictures of ray
orders $4-j$, for $j=1,2,3$. The Wilson pictures have ray orders
$3-j$. Their sums therefore have the required ray orders $3-j$. Prices
and exact recording identities add. The quadratic normalization and the
cubic vanishing of the curvature part give $U_m(0)=0$ and
$DU_m(0)=0$. Apply Proposition~\ref{prop:ex-three-level-conversion}.
Its finite constants depend only on the listed bounds and the finite
color, phase, and stitch sets, so the common width and price are uniform
in the path, period, and selected action.
\end{proof}

\paragraph{Scope.}
This section proves the uniform atom, Ward, and Wilson-subtracted
coefficient estimates, together with the finite division and the exact
stitched response identity. The dependence of downstream applications on
the displayed metric, floor, and sworn hypotheses remains as stated.
